\cotefile{32}
\scanpage{37}
Pour tout carré cartésien de schémas
\diagram{\begin{tikzcd}[ampersand replacement=\&]
 X' \arrow[r,"g'"] \arrow[d,"f'"] \& X \arrow[d,"f"]\\
 S' \arrow[r,"g"] \& S
\end{tikzcd}}
avec $f$ compactifiable, \deleted{\unsure{$\mathcal C$|$\ell$|$\mathcal E$}} et tout faisceau \inserted{d'anneaux} de torsion $\mathcal A$ sur $S$, l'isomorphisme de changement de base $\overline{\underline{\mathrm{XVII}}}$ 5.2
\margin{\unsure{$p$}}
\[
 g^* Rf_! \xrightarrow{\sim} Rf'_! g'^*
\]
entre foncteurs de $D(X,f^*\mathcal A)$ dans $D(S',g^*\mathcal A)$ se transpose en un isomorphisme
\[
 Rg'_* Rf'{}^! \xrightarrow{\sim} Rf^! Rg_* \tag*{(3.1.6.3)}
\]
entre foncteurs de $D^+(S',g^*\mathcal A)$ dans $D^+(X,f^*\mathcal A)$.

On obtient ainsi quatre ``\uline{flèches de changement de base}'' :
\diagram{\begin{tikzcd}[ampersand replacement=\&,column sep=large,row sep=large,
 execute at end picture={\draw[rounded corners,->] ([xshift=5mm,yshift=4mm]\tikzcdmatrixname-2-1.north east) -- ++(0,-3mm) -| ([xshift=-3mm]\tikzcdmatrixname-2-1.west) |- ([xshift=4mm,yshift=-3mm]\tikzcdmatrixname-3-1.south east) -- ++(0,-3mm);
 \draw[->] (\tikzcdmatrixname-1-1.south west) to[bend right=25] ([xshift=-6mm,yshift=-2mm]\tikzcdmatrixname-2-1.west);}]
 \tikz[baseline=(n.base)]{\node(n){$(3.1.6.4)$};\draw (n.center) ellipse (9mm and 3mm);} \& {} \& {} \& {}\\
 \deleted{\unsure{R}}g^*Rf_* \arrow[r] \& Rf'_*g'^* \& g^*Rf_! \arrow[r,"\sim", "\deleted{\leftarrow}"' near start] \& Rf'_!g'^*\\
 Rg^!Rf_* \& Rf'_*Rg'{}^! \arrow[l,"\sim"',"\deleted{\rightarrow}" near start] \& Rg^!Rf_! \& Rf'_!Rg'{}^! \arrow[l]
\end{tikzcd}}
Les flèches de la diagonale principale sont autoadjointes, et définies pour tout carré commutatif dès que leurs sources et buts ont un sens, \inserted{et en tous cas quand on considère les foncteurs envisagés comme définis sur \inserted{des catégories $D^+$.}} les isomorphismes de la seconde diagonale sont adjoints l'un de l'autre, et définis pour un carré \deleted{commutatif}\inserted{cartésien} de \deleted{\unsure{petits}} schémas, avec $f$ (resp $g$) compactifiable.

\inserted{Chacune de} ces flèches vérifie\deleted{\unsure{nt}} \deleted{des}\inserted{deux} compatibilités du type $\overline{\underline{\mathrm{XII}}}$ 4.4 pour $f$ ou $g$ composé de deux morphismes.
% slanted
\margin{Peut-on dire la raison?}
\margin{$g'^*Rf^!\xrightarrow{\sim}Rf'{}^!g^*$}
% slanted
\margin{\unsure{défini|dû} sans hyp. de finité}
% faint pencil at the foot, with a line between the two groups
\margin{(\unsure{est-ce que Verdier} \unsure{utilise} \unsure{pour calculer} $Rf^!$ \unsure{son} \illegible{1} $\longleftarrow$)}
\margin{(\unsure{manque} \illegible{1} : $\underline g$ loc. acyclique (p. ex. \uline{lisse}))}
% archive sheet number: 36
